% SHARD-ID: AQM-22-PRIME-FIELD-ARITHMETIC-FIELDS-QUDIT-PAULIS
% SHARD-TITLE: Prime-Field Arithmetic Fields and Qudit Paulis
% SHARD-SUMMARY: Separates the one-point spectrum of a prime field from its underlying set of elements.
% SHARD-SUMMARY: Proves that the all-map arithmetic field on \(X=\mathbb F_p\) gives the \(p\)-qudit Pauli group.
% SHARD-SUMMARY: Records the \(p=2\) phase edge case and the distinction between ambient Pauli groups and stabilizer choices.
% SHARD-KEYWORDS: arithmetic quantum field, finite field, prime field, qudit, Pauli group, Weyl-Heisenberg, stabilizer

\section{Prime-Field Arithmetic Fields and Qudit Paulis}
\label{sec:prime-field-arithmetic-fields-qudit-paulis}

\subsection{Status and Source Anchors}

This shard stays below schemes and sheaves.  The only commutative rings are
prime fields \(\mathbb F_p\).  The local Pauli/Weyl sources are Ashikhmin--Knill,
\path{references/symplectic_qecc/ashikhmin_knill_0005008_source/n9.tex}, lines 249--313; Gottesman's thesis, \path{references/symplectic_qecc/gottesman_9705052_source/Thesis.tex}, lines 1627--1681; and Bombin--Martin--Delgado, \path{references/toric_code/bombin_martin_delgado_0605094_source/HomologicalErrorCorrection6.tex}, lines 1208--1259.

\paragraph{Source-sign bridge.}
The local convention is \path{CONVENTIONS.md}, convention (t): \(W(x,z)=T_xR_z\) with \(T_x|s\rangle=|s+x\rangle\).  The cited source formulas use the opposite commutator exponent if their displayed shift exponent is identified directly with this local \(X\)-coordinate; see Ashikhmin--Knill \path{references/symplectic_qecc/ashikhmin_knill_0005008_source/n9.tex}, lines 265--268 and 307--313; Gottesman \path{references/symplectic_qecc/gottesman_9705052_source/Thesis.tex}, lines 1677--1680; and Bombin--Martin--Delgado \path{references/toric_code/bombin_martin_delgado_0605094_source/HomologicalErrorCorrection6.tex}, lines 1247--1254.  This shard uses the conversion
\(x_{\rm source}=-x_{\rm local}\), \(z_{\rm source}=z_{\rm local}\), so the source exponent \(xz'-zx'\) becomes the local exponent \(zx'-xz'\).  Agreement with the cited Pauli group is therefore after inverting the source shift generator, not identical displayed sign conventions.

\subsection{Lamport Dependency Skeleton}

\[
\begin{array}{rcl}
D1&:&p\text{ is prime and }R=\mathbb F_p.\\
L1&:&R\text{ has one prime ideal and }\operatorname{Frac}(R)=R.\\
D2&:&X_{\rm sp}=\operatorname{Spec}R,\quad X_{\rm el}=R\text{ as a finite set.}\\
D3&:&V=R_X\oplus R_Z,\quad
       \omega_V((x,z),(x',z'))=zx'-xz'.\\
D4&:&E_{\rm el}=\operatorname{Map}(X_{\rm el},V).\\
D5&:&\Omega((\xi,\zeta),(\xi',\zeta'))=
       \sum_{u\in X_{\rm el}}(\zeta(u)\xi'(u)-\xi(u)\zeta'(u)).\\
D6&:&\mathcal H_X=\ell^2(R^{X_{\rm el}}),\quad W(\xi,\zeta)=T_\xi R_\zeta.\\
L2&:&(V,\omega_V)\text{ and }(E_{\rm el},\Omega)\text{ are symplectic.}\\
L3&:&W(\xi,\zeta)W(\xi',\zeta')
       =\chi(\zeta\cdot\xi')W(\xi+\xi',\zeta+\zeta').\\
T1&:&\langle\chi(c)I,W(E_{\rm el})\rangle
       \text{ is the standard odd-}p\text{ qudit Pauli group on }p
       \text{ qudits after the source-sign bridge.}\\
T2&:&\text{Stabilizers require an extra isotropic label subspace and phases.}\\
L4&:&p=2\text{ has the same projective labels but needs the usual }\mu_4
       \text{ Pauli phases for Hermitian }Y.
\end{array}
\]

\subsection{Prime Field Points L1--D2}

\paragraph{Lemma \(L1\).}
The ring \(R=\mathbb F_p\) has exactly one prime ideal, namely \((0)\), and
\(\operatorname{Frac}(R)=R\).

\emph{Proof.}
A field has only two ideals, \((0)\) and the whole field.  The whole field is
not prime because prime ideals are proper.  The zero ideal is prime because
the quotient \(R/(0)=R\) is an integral domain.  Thus
\(\operatorname{Spec}R=\{(0)\}\).  The fraction field consists of symbols
\(a/b\) with \(b\ne0\), and the map \(a/b\mapsto ab^{-1}\) is a well-defined
field isomorphism \(\operatorname{Frac}(R)\cong R\).

There are therefore two concrete point sets in play.  The spectrum
\[
  X_{\rm sp}=\operatorname{Spec}\mathbb F_p
\]
has one point.  The underlying set of field elements
\[
  X_{\rm el}=\mathbb F_p
\]
has \(p\) points.  The one-point choice gives one qudit.  The \(p\)-point
choice is the one that can give the \(p\)-qudit Pauli group.

\subsection{The All-Map Arithmetic Field D3--L2}

Let
\[
  V=R_X\oplus R_Z
\]
with symplectic form
\[
  \omega_V((x,z),(x',z'))=zx'-xz'.
\]
For \(X=X_{\rm el}\), define
\[
  E_{\rm el}=\operatorname{Map}(X,V)
  \cong R^X\oplus R^X.
\]
Write an element of \(E_{\rm el}\) as \((\xi,\zeta)\), where
\(\xi,\zeta:X\to R\).  Define
\[
  \Omega((\xi,\zeta),(\xi',\zeta'))
  =
  \sum_{u\in X}\bigl(\zeta(u)\xi'(u)-\xi(u)\zeta'(u)\bigr).
\]

\paragraph{Lemma \(L2\).}
\((V,\omega_V)\) and \((E_{\rm el},\Omega)\) are symplectic \(R\)-vector
spaces.  Moreover \(E_{\rm el}\cong R^p\oplus R^p\).

\emph{Proof.}
The form \(\omega_V\) is bilinear.  It is alternating because
\[
  \omega_V((x,z),(x,z))=zx-xz=0.
\]
If \((x,z)\ne0\), then \(\omega_V((x,z),(1,0))=z\) and
\(\omega_V((x,z),(0,1))=-x\), so some pairing is nonzero.  Thus \(V\) is
symplectic.

Since \(X\) has \(p\) elements, choosing any ordering of \(X\) identifies
\(R^X\) with \(R^p\).  Hence \(E_{\rm el}\cong R^p\oplus R^p\).  The form
\(\Omega\) is the direct sum of the pointwise symplectic forms.  If
\((\xi,\zeta)\ne0\), choose \(u_0\) with \((\xi(u_0),\zeta(u_0))\ne0\) and use
the previous paragraph to choose a value at \(u_0\) pairing nontrivially,
with all other values zero.  Thus \(\Omega\) is nondegenerate and alternating.

\paragraph{Map versus Hom.}
The \(p\)-qudit statement uses all maps.  If instead one chooses the
structure-preserving space
\[
  \operatorname{Hom}_R(R,V),
\]
where \(R\) is viewed as a one-dimensional \(R\)-vector space, then evaluation
at \(1\) gives \(\operatorname{Hom}_R(R,V)\cong V\).  That smaller choice is
the one-qudit label space, not the \(p\)-qudit label space.

\subsection{Weyl Operators L3}

Let \(\chi(t)=\exp(2\pi i\,\tilde t/p)\), where \(\tilde t\in\{0,\ldots,p-1\}\)
represents \(t\in R\).  On the Hilbert space
\[
  \mathcal H_X=\ell^2(R^X),
\]
with basis \(\{|s\rangle:s:X\to R\}\), define
\[
  T_\xi|s\rangle=|s+\xi\rangle,\qquad
  R_\zeta|s\rangle=\chi\!\left(\sum_{u\in X}\zeta(u)s(u)\right)|s\rangle,
\]
and
\[
  W(\xi,\zeta)=T_\xi R_\zeta.
\]

\paragraph{Lemma \(L3\).}
For \(e=(\xi,\zeta)\) and \(f=(\xi',\zeta')\),
\[
  W(e)W(f)=\chi(\zeta\cdot\xi')W(e+f),
\]
where \(\zeta\cdot\xi'=\sum_u\zeta(u)\xi'(u)\).  Hence
\[
  W(e)W(f)=\chi(\Omega(e,f))W(f)W(e).
\]

\emph{Proof.}
The translation operators add and the phase operators add:
\[
  T_\xi T_{\xi'}=T_{\xi+\xi'},\qquad
  R_\zeta R_{\zeta'}=R_{\zeta+\zeta'}.
\]
Also
\[
\begin{aligned}
  R_\zeta T_{\xi'}|s\rangle
  &=R_\zeta|s+\xi'\rangle\\
  &=\chi(\zeta\cdot s+\zeta\cdot\xi')|s+\xi'\rangle\\
  &=\chi(\zeta\cdot\xi')T_{\xi'}R_\zeta|s\rangle .
\end{aligned}
\]
Therefore
\[
  T_\xi R_\zeta T_{\xi'}R_{\zeta'}
  =
  \chi(\zeta\cdot\xi')T_{\xi+\xi'}R_{\zeta+\zeta'}.
\]
Swapping \(e\) and \(f\) gives the commutator factor
\[
  \chi(\zeta\cdot\xi'-\zeta'\cdot\xi)=\chi(\Omega(e,f)).
\]

\subsection{Identification with Qudit Paulis T1--T2}

For each \(u\in X\), let \(\delta_u\in R^X\) be the delta function.  Define
single-site operators
\[
  X_u=W(\delta_u,0),\qquad Z_u=W(0,\delta_u).
\]
Let \(\mathcal P_p^\chi(X)\) be the group generated by all \(X_u,Z_u\) and the
central character phases \(\chi(c)I\), \(c\in R\).  For odd \(p\), this is the
standard prime-dimensional qudit Pauli group after the source-sign conversion
above: one \(p\)-dimensional qudit for each \(u\in X\), shifts generated by
\(X_u\), and phases generated by \(Z_u\).  For \(p=2\), it is the unphased Weyl
group whose projective labels are the usual qubit Pauli labels; the
conventional Hermitian \(Y\) convention is recorded in \(L4\).

\paragraph{Theorem \(T1\).}
For \(X=\mathbb F_p\) as a set of elements,
\[
  \langle \chi(c)I,\ W(e):c\in R,\ e\in E_{\rm el}\rangle
  =
  \mathcal P_p^\chi(X).
\]
Thus, for odd \(p\), the all-map arithmetic quantum field on
\(X=\mathbb F_p\) with target
\(\operatorname{Frac}(\mathbb F_p)^2=\mathbb F_p^2\) gives the ambient
Gottesman/Ashikhmin--Knill \(p\)-qudit Pauli group for \(p\) qudits, after
inverting the source shift generator as stated in the source-sign bridge.  For
\(p=2\), the projective statement is the same and the full conventional Pauli
phase convention is \(L4\).

\emph{Proof.}
The operators \(W(\delta_u,0)\) and \(W(0,\delta_u)\) are among the arithmetic
field Weyl operators, so \(\mathcal P_p^\chi(X)\) is contained in the generated
arithmetic Weyl group.  Conversely, for any \((\xi,\zeta)\),
\[
  T_\xi=\prod_{u\in X}X_u^{\xi(u)},\qquad
  R_\zeta=\prod_{u\in X}Z_u^{\zeta(u)},
\]
with any fixed site order, because different \(X_u\)'s commute and different
\(Z_u\)'s commute.  Thus \(W(\xi,\zeta)=T_\xi R_\zeta\) lies in
\(\mathcal P_p^\chi(X)\).  The two generated groups are equal.

\paragraph{Theorem \(T2\).}
The arithmetic field datum gives the ambient Pauli/Weyl group.  A stabilizer
code is extra data: an isotropic label subspace \(L\leq E_{\rm el}\), plus a
compatible choice of phases/eigenvalue character.

\emph{Proof.}
If \(S\leq\mathcal P_p^\chi(X)\) is Abelian, its image \(L\) in the projective label
space \(E_{\rm el}\) satisfies \(\Omega(L,L)=0\) by Lemma \(L3\).  Conversely,
if \(p\) is odd and \(L\leq E_{\rm el}\) is isotropic, define
\[
  \widetilde W(\xi,\zeta)=
  \chi\!\left(\frac12\,\zeta\cdot\xi\right)W(\xi,\zeta).
\]
Using Lemma \(L3\),
\[
  \widetilde W(e)\widetilde W(f)
  =
  \chi\!\left(\frac12\Omega(e,f)\right)\widetilde W(e+f).
\]
On an isotropic \(L\), this becomes
\(\widetilde W(e)\widetilde W(f)=\widetilde W(e+f)\).  Multiplying this lift by
a character \(L\to\mathbb C^\times\) changes the stabilizer eigenvalues.  Thus
the arithmetic field supplies the ambient labels and commutation form; the
choice of stabilizer is a later choice of isotropic subspace and phases.

\subsection{The p=2 Edge Case L4}

For \(p=2\), the same label space \(E_{\rm el}\), the same Hilbert space, and
the same commutator criterion
\[
  W(e)W(f)=(-1)^{\Omega(e,f)}W(f)W(e)
\]
remain valid.  The edge is the phase convention.  There is no element
\(1/2\in\mathbb F_2\), and the unphased operator \(W(1,1)=XZ\) is not the
Hermitian Pauli \(Y\).  The standard qubit Pauli group uses fourth roots of
unity and the relabeling
\[
  \widehat W(\xi,\zeta)=
  i^{\sum_u \xi(u)\zeta(u)}W(\xi,\zeta),
\]
so \(\widehat W(\delta_u,\delta_u)=Y_u\).  Hence \(p=2\) has the same
projective arithmetic field labels and the same isotropy criterion, but the
full conventional Pauli group has center \(\mu_4\), not only the character
phases \(\mu_2\).

\subsection{Conclusion}

The conjecture is correct with one precision.  If \(X=\operatorname{Spec}
\mathbb F_p\), the construction gives one qudit.  If \(X=\mathbb F_p\) is the
underlying \(p\)-element set and \(E=\operatorname{Map}(X,\mathbb F_p^2)\), the
Weyl--Heisenberg arithmetic field gives exactly the standard odd-\(p\) qudit
Pauli group, and for \(p=2\) gives the same projective Pauli labels with the
usual fourth-root phase adjustment.  In every prime case the projective labels
are
\[
  E\cong\mathbb F_p^p\oplus\mathbb F_p^p.
\]
Actual stabilizer codes are not automatic fields; they are isotropic
subspaces of this field label space together with phase/eigenvalue choices.
